\section{Initialization and spin readout}
\label{sec:initialization-readout}

Initialization prepares a known logical input; readout converts the
logical basis populations into distinguishable electrical signals.
The readout schemes considered here first convert the spin
distinction into different charge configurations or tunneling
records. A nearby charge sensor detects these through their
electric effect on its current (Appendix~\ref{sec:measurement-charge-sensor}).

\subsection{Single-spin qubits}
\label{sec:single-spin-spam}
\label{sec:spin-initialization}

The initialization target is $\ket0=\ket{\downarrow}$, and readout
distinguishes $\ket{\downarrow}$ from $\ket{\uparrow}$.
We suppress the common orbital--valley state in the notation.
The lower spin state can be prepared either by relaxation after
loading or by selecting it during loading.

\paragraph{Initialization by relaxation.}
First load one electron from a reservoir. Its spin may initially
occupy either Zeeman level. Raising the reservoir barrier then
suppresses electron exchange and keeps the dot occupied.
The spin remains coupled to the thermal environment, so an excited
$\ket{\uparrow}$ can release energy and relax to the lower
$\ket{\downarrow}$ state. Waiting for this dissipative process
prepares the input state; a unitary rotation alone cannot reset an
arbitrary unknown spin.

The achievable ground-state population and required waiting time
follow from the relaxation dynamics. Let $p_\uparrow(t)$ be the
probability of occupying the upper spin level.
For a bath at temperature $T$, its equilibrium value is
$p_\uparrow^{\mathrm{eq}}=[1+\exp(E_Z/k_BT)]^{-1}$, where $k_B$
is Boltzmann's constant. A simple relaxation model is
\begin{equation}
  \dot p_\uparrow=-\frac{p_\uparrow-p_\uparrow^{\mathrm{eq}}}{T_1},
  \qquad
  p_\uparrow(t)=p_\uparrow^{\mathrm{eq}}
       +[p_\uparrow(0)-p_\uparrow^{\mathrm{eq}}]e^{-t/T_1}.
  \label{eq:thermal-spin-initialization}
\end{equation}
Here $T_1$ is the longitudinal relaxation time.
The initial-state dependence decays on the scale $T_1$, but the
remaining excited-state population is set by the bath temperature.
High-probability ground-state preparation therefore requires both
a wait much longer than $T_1$ and $E_Z\gg k_BT$.
Use the measured $T_1$ and long-time excited-state population at the
holding voltages to choose the waiting time and determine the
attainable preparation probability
(Appendix~\ref{sec:measurement-relaxation}). The loading and isolation
settings are specified in
Appendix~\ref{sec:calibration-single-spin-initialization}.

\paragraph{Spin-selective loading.}
The second method prepares the lower spin state as the electron
enters the dot, without relying on subsequent spin relaxation.
Start from an empty dot. The energy required to add the first
electron is $\mu_\sigma=E(1,\sigma)-E(0)$, where $E(0)$ and
$E(1,\sigma)$ are the empty-dot and one-electron energies.
Its spin dependence comes from the Zeeman energy of the final state.
If $\epsilon_d$ is the addition energy without the Zeeman term,
\begin{equation}
  \mu_\downarrow=\epsilon_d-\frac{E_Z}{2},\qquad
  \mu_\uparrow=\epsilon_d+\frac{E_Z}{2}.
  \label{eq:spin-resolved-addition-energies}
\end{equation}
Thus the two loading thresholds differ by $E_Z$ even for the first
electron, without an electron--electron interaction.
Gate voltages shift their common position $\epsilon_d$ relative to
the reservoir chemical potential $\mu_{\mathrm R}$.
To allow down-spin loading while suppressing up-spin loading, set
\begin{equation}
  \mu_\downarrow<\mu_{\mathrm R}<\mu_\uparrow.
  \label{eq:spin-selective-window}
\end{equation}
The spin-resolved addition thresholds must be measured to select
the plunger voltage that satisfies this window
(Appendix~\ref{sec:calibration-single-spin-readout}).
At low temperature, reservoir states at $\mu_\downarrow$ are
occupied and can supply an electron, while states at
$\mu_\uparrow$ are mostly empty. Charging energy must also prevent
a second electron from entering.

The level alignment determines which process is allowed; the tunnel
rates determine how long to keep the loading window open.
For weak spin-conserving tunneling, let $\Gamma_\sigma$ be the bare
tunnel rate and $T_e$ the reservoir electron temperature.
The Fermi occupation factor gives the conditional rates
\begin{equation}
  \begin{aligned}
    f_{\mathrm R}(E)&=\frac{1}{1+\exp[(E-\mu_{\mathrm R})/(k_BT_e)]},\\
    \Gamma_{\mathrm{in},\sigma}&=\Gamma_\sigma f_{\mathrm R}(\mu_\sigma),
    \qquad
    \Gamma_{\mathrm{out},\sigma}
       =\Gamma_\sigma[1-f_{\mathrm R}(\mu_\sigma)].
  \end{aligned}
  \label{eq:reservoir-spin-rates}
\end{equation}
The factors $f_{\mathrm R}$ and $1-f_{\mathrm R}$ require an
occupied source state and an empty destination state, respectively.
Appendix~\ref{sec:amplitudes-to-master-equation} derives these
reservoir rates and their waiting-time law from tunneling amplitudes.
Both margins in Eq.~\eqref{eq:spin-selective-window} should exceed
$k_BT_e$ and tunnel broadening so that loading and unloading remain
spin-selective.

Neglecting other events, an initially empty dot remains empty with
probability $e^{-\Gamma_{\mathrm{in},\downarrow}t}$.
Measure the loading rate at the chosen reservoir-barrier setting
(Appendix~\ref{sec:calibration-single-spin-initialization}) and use
it to choose a loading interval much longer than
$\Gamma_{\mathrm{in},\downarrow}^{-1}$. Then raise the reservoir
barrier to retain the electron during the gate sequence.
Spin-selective loading has been demonstrated in
$\Si/\Si\Ge$ dots~\cite{Simmons2011}.
Excited orbitals and energy-dependent barriers can modify the rates.

\paragraph{Readout by energy-selective tunneling.}
\label{sec:energy-selective-readout}

The readout goal is to determine whether the occupied dot contains
an up-spin or down-spin electron after the gates.
We convert that distinction into a charge transient by allowing
only the upper spin level to empty into the reservoir.

Reopen the reservoir barrier and tune the addition energies to
Eq.~\eqref{eq:spin-selective-window}.
An up-spin electron can tunnel into an empty reservoir state;
a down-spin electron remains because the corresponding states
are occupied. If the dot empties, it can reload a down-spin electron.
The two initial spin states therefore give different charge records:
\begin{equation}
  \begin{aligned}
    \ket{\uparrow}:&\quad N=1\longrightarrow0\longrightarrow1,\\
    \ket{\downarrow}:&\quad N=1\quad\text{throughout the readout interval}.
  \end{aligned}
  \label{eq:energy-readout-records}
\end{equation}
A sensor dot or quantum point contact detects the temporary empty
state through a current or radio-frequency response~\cite{Elzerman2004}.
The evidence for an initial up spin is the $1\to0\to1$ transient,
not the final occupation, which can be one for either spin.

The electron must escape before it relaxes: a spin that relaxes first
leaves no up-spin charge record. To quantify this competition, take
constant escape and relaxation rates
$\Gamma_{\mathrm{out},\uparrow}$ and $T_1^{-1}$
during a low-temperature readout interval $\tau_{\mathrm r}$.
The probability of neither event by time $t$ is
$e^{-(\Gamma_{\mathrm{out},\uparrow}+T_1^{-1})t}$.
Multiplying by the escape rate and integrating gives the probability
of producing the desired charge transient:
\begin{equation}
  P_{\mathrm{escape}|\uparrow}(\tau_{\mathrm r})
    =\frac{\Gamma_{\mathrm{out},\uparrow}}
           {\Gamma_{\mathrm{out},\uparrow}+T_1^{-1}}
      \left[1-e^{-(\Gamma_{\mathrm{out},\uparrow}+T_1^{-1})\tau_{\mathrm r}}\right].
  \label{eq:readout-competing-rates}
\end{equation}
Even when escape occurs, the detector must resolve the empty interval.
Fast reloading or sensor noise can hide it, while thermal down-spin
escape can produce a false up-spin signal.
Thus readout fidelity depends on the charge-conversion probability
above and on the sensor response. The measured escape rates and
$T_1$ at the readout voltages determine the reservoir-barrier setting
and readout interval: up-spin escape should be likely before
relaxation, while down-spin escape remains unlikely
(Appendix~\ref{sec:calibration-readout-duration-choice}).

For $a\ket{\downarrow}+b\ket{\uparrow}$, ideal readout gives the
up-spin outcome with probability $|b|^2$. Measurements after
calibrated rotations on repeated preparations access other spin
components. Escape and reloading make this protocol destructive.
Appendix~\ref{sec:calibration-single-spin-readout} specifies the
conversion window, and Appendix~\ref{sec:calibration-readout-duration-choice}
derives the competing-rate probability and duration constraints.

\subsection{Singlet--triplet qubits}
\label{sec:st-spam}
\label{sec:pauli-blockade-readout}

Here initialization prepares $\ket{0_{\mathrm{ST}}}=\ket{S_{11}}$,
and readout distinguishes it from $\ket{1_{\mathrm{ST}}}=\ket{T_{0,11}}$.
Both tasks use the singlet's access to double occupation.

\paragraph{Initialization.}
A convenient route first puts both electrons on one dot, where the
lowest orbital supports a unique spin singlet. Tune to a region
where $\ket{D_{02}}$ is the nondegenerate ground state, load two
electrons, and allow relaxation. Then sweep across the singlet
charge anticrossing to separate them into $(1,1)$ while preserving
the singlet spin state. The sweep must follow the charge anticrossing adiabatically.
The measured tunnel coupling and detuning response determine the
allowed voltage-ramp speed through the anticrossing. Once exchange
becomes small, use the measured Zeeman difference to limit the time
available for singlet--triplet rotation, or compensate that rotation
(Appendix~\ref{sec:calibration-st-initialization}).
Accessible orbital or valley excitations can spoil this preparation.

\paragraph{Readout by Pauli spin blockade.}
The goal is to map the singlet and triplet to different charge states. Start in $(1,1)$ and
pulse toward $(0,2)$.
The singlet couples to $\ket{D_{02}}$ through the tunnel matrix
element in Eq.~\eqref{eq:tunnel-singlet-matrix-elements}.
A doubly occupied triplet needs an excited orbital--valley state;
keeping that state outside the readout window blocks triplet transfer.
Choose the readout detuning from the measured singlet and triplet
charge-transfer thresholds so that only the singlet transition is
allowed (Appendix~\ref{sec:calibration-st-readout}).
With negligible spin mixing and charge-transfer error,
\begin{equation}
  \ket{S_{11}}\longrightarrow\ket{D_{02}},\qquad
  \ket{T_{11}}\longrightarrow\ket{T_{11}},
  \label{eq:pauli-readout-map}
\end{equation}
Here $\ket{T_{11}}$ is a triplet within the blocked window.
A charge sensor distinguishes the resulting $(0,2)$ and $(1,1)$
occupations, implementing singlet--triplet readout through
Pauli spin blockade~\cite{Wu2014}.
Because the two configurations have the same total electron number,
the sensor must have different sensitivity to charge on the two
dots (Appendix~\ref{sec:measurement-charge-sensor}).
This pulse deliberately permits real charge transfer, so the
weak-admixture approximation used to derive exchange does not
describe the transfer itself.

The measurement tests whether the input has a singlet component.
For an initial two-spin density operator $\rho$, the probabilities
of the two charge outcomes are
\begin{equation}
  p_{02}=\operatorname{Tr}(\rho P_S),\qquad
  p_{11}=\operatorname{Tr}[\rho(I-P_S)],\qquad
  P_S=\ket{S_{11}}\bra{S_{11}}.
  \label{eq:pauli-readout-probabilities}
\end{equation}
Within the singlet--triplet encoding, these probabilities are the
populations of $\ket{0_{\mathrm{ST}}}$ and $\ket{1_{\mathrm{ST}}}$.
A blocked outcome outside that encoding cannot distinguish $T_0$
from $T_\pm$. Relaxation or field-gradient-induced mixing can turn
a blocked state into one that transfers, so the sensor must resolve
the charge outcome before those processes erase the distinction.
Appendix~\ref{sec:calibration-st-readout} develops the readout
window and the conditions for resolving the two charge outcomes.


\subsection{Exchange-only qubits}
\label{sec:exchange-only-spam}

For the fixed-$m$ encoding in Eq.~\eqref{eq:exchange-only-basis},
initialization prepares $\ket{0_{\mathrm E}}$, and readout must
distinguish it from $\ket{1_{\mathrm E}}$.
The useful distinction is the spin of pair 1--2: a singlet in
$\ket{0_{\mathrm E}}$ and a triplet in $\ket{1_{\mathrm E}}$.

\paragraph{Initialization.}
Let $\ket{S_{12}}=(\ket{\uparrow\downarrow}
-\ket{\downarrow\uparrow})/\sqrt2$ denote the singlet of spins 1 and 2.
The chosen logical zero factorizes as
\begin{equation}
  \ket{0_{\mathrm E}}
    =\ket{S_{12}}\otimes\ket{\uparrow}_3
    =\frac{\ket{\uparrow\downarrow\uparrow}
                  -\ket{\downarrow\uparrow\uparrow}}{\sqrt2}.
  \label{eq:eo-initial-state}
\end{equation}
We therefore prepare a singlet on dots 1--2 and an up spin on dot 3,
with tunneling between dots 2 and 3 suppressed during preparation.
Use the double-occupation preparation described above for the pair:
load its singlet in $(0,2)$, then separate the electrons without
mixing singlet and triplet. Keeping one electron on dot 3, the
charge sequence is $(0,2,1)\to(1,1,1)$.
The singlet-loading and charge-transfer steps have the same
temperature, excited-state, and ramp constraints as for the
singlet--triplet qubit.

The up spin on dot 3 is an additional requirement of our choice
$m=+1/2$. With the Zeeman convention used here, relaxation prepares
the opposite spin. If individually addressable transverse control
is available, one route is to load $\ket{\downarrow}_3$ as in
Section~\ref{sec:single-spin-spam}, then apply a resonant $\pi$ pulse
to spin 3 before the encoded gates. Its frequency and duration are
set by the measured resonance and rotation rate of spin 3
(Appendix~\ref{sec:measurement-rabi}). This preparation uses a control
beyond exchange; the subsequent logical rotations use exchange alone.
An exchange pulse cannot set $m$, because it conserves
the total spin projection along $z$.

Polarizing the third spin is not required by the broader
exchange-only subsystem encoding used in silicon
experiments~\cite{Weinstein2023}.
Preparing only the pair singlet gives
$\rho_{\mathrm{init}}=\ket{S_{12}}\bra{S_{12}}\otimes\rho_3$,
where $\rho_3$ is the third spin's density operator.
For example, an unpolarized $\rho_3=I_3/2$ populates both
$m=\pm1/2$ sectors. Both correspond to logical zero because pair
1--2 is a singlet; the value of $m$ is left unspecified.
This is a larger encoding than the fixed-$m$ subspace used above.
Using it for two-qubit gates requires a sequence that implements
the same logical gate for either value of each qubit's $m$, as in
the Fong--Wandzura construction~\cite{FongWandzura2011,Weinstein2023}.
Appendix~\ref{sec:calibration-eo-initialization} gives the preparation
steps and distinguishes fixed-$m$ from subsystem initialization.

\paragraph{Readout of pair 1--2.}
To measure the logical basis, suppress $J_{23}$ and pulse dots 1--2
toward double occupation while dot 3 remains singly occupied.
A pair singlet can enter the lowest $(0,2)$ state; a pair triplet
is blocked when the excited doubly occupied triplets lie outside
the readout window. The sensor therefore distinguishes $(0,2,1)$
from $(1,1,1)$, without measuring spin 3.
Pair singlet preparation and Pauli-blockade readout are used in
silicon triple-dot qubits~\cite{Blumoff2022}.

To see why this measures the encoded basis, write the singlet
projector on the three-spin space as
\begin{equation}
  Q_S=\ket{S_{12}}\bra{S_{12}}\otimes I_3
       =\frac14 I-\mathbf{s}_1\cdot\mathbf{s}_2.
  \label{eq:eo-readout-projector}
\end{equation}
Here $I_3$ acts on spin 3 and $I$ on all three spins.
Pair exchange has eigenvalue $-3/4$ on the pair singlet and $1/4$
on any pair triplet, so
\begin{equation}
  Q_S\ket{0_{\mathrm E}}=\ket{0_{\mathrm E}},\qquad
  Q_S\ket{1_{\mathrm E}}=0,\qquad
  \left.P_{\mathrm E}Q_SP_{\mathrm E}\right|_{\mathrm E}
        =\frac{I_{\mathrm E}+Z_{\mathrm E}}{2},
  \label{eq:eo-logical-readout}
\end{equation}
where $P_{\mathrm E}$ projects onto the two-dimensional encoded
space, $I_{\mathrm E}$ is its identity, and
$Z_{\mathrm E}=\ket{0_{\mathrm E}}\bra{0_{\mathrm E}}
-\ket{1_{\mathrm E}}\bra{1_{\mathrm E}}$.
Consequently, ideal charge conversion for
$\ket{\psi}=a\ket{0_{\mathrm E}}+b\ket{1_{\mathrm E}}$ gives
\begin{equation}
  p_{021}=|a|^2,\qquad p_{111}=|b|^2.
  \label{eq:eo-readout-probabilities}
\end{equation}
The subscripts label the three-dot occupations.
A preceding encoded rotation changes the measured logical axis.

This assignment assumes that the input is in the computational
space. States with total spin $S=3/2$ also have a triplet on pair
1--2 and produce the blocked outcome; pair readout alone cannot
distinguish that leakage from logical one.
It also does not determine $m$. Thus the same singlet test works
for the subsystem encoding with an unpolarized third spin, but
does not certify preparation in our fixed-$m$ subspace.
Residual $J_{23}$, magnetic-field gradients, or relaxation during
conversion can change the pair spin before it is resolved.
Their timescales must be long compared with conversion and sensing.
Appendix~\ref{sec:calibration-eo-readout} develops this pair-readout
protocol and its limitations.

\subsection{From repeated records to logical-state probabilities}
\label{sec:platform-readout-probabilities}

Each readout returns one electrical record, such as an empty-dot
transient or a charge configuration after Pauli-blockade conversion.
The conversion rules above assign that record a logical outcome.
For an input density operator $\rho$ within the chosen qubit space,
ideal logical-basis measurement has probabilities
\begin{equation}
 p_0=\operatorname{Tr}(\rho\ket0\bra0),\qquad
 p_1=1-p_0,\qquad
 \widehat p_0=\frac{n_0}{M}.
 \label{eq:platform-repeated-readout}
\end{equation}
Here $\ket0,\ket1$ denote the logical basis of the encoding under
study, and $n_0$ is the number of zero outcomes in $M$ repetitions.
Each repetition prepares the input again and applies the same
control sequence before readout. A single outcome is random;
the fraction $n_0/M$ estimates its probability. For independent
repetitions with fixed $p_0$, its variance is $p_0(1-p_0)/M$.
Readout errors change the relation between the observed frequency
and the logical population. If logical 0 is reported as 1 with
probability $\epsilon_0$, and logical 1 as 0 with probability
$\epsilon_1$, fixed classical assignment errors give
\begin{equation}
 p_{1,\mathrm{obs}}
   =\epsilon_0+(1-\epsilon_0-\epsilon_1)p_1.
 \label{eq:platform-readout-assignment}
\end{equation}
Readout thus introduces an offset and reduces contrast. Recovering
$p_1$ requires nonzero contrast and independently established error
probabilities. Appendix~\ref{sec:calibration-record-classification}
derives this assignment map and its inverse in
Eq.~\eqref{eq:calibration-assignment-map}.

Logical populations alone do not determine a quantum state.
For example, $\ket+= (\ket0+\ket1)/\sqrt2$ and the mixed state
$I/2$ both give $p_0=p_1=1/2$. Their difference appears if a known
logical rotation $U_{\mathrm a}$ precedes the same readout.
The probability after an analysis rotation is then
\begin{equation}
 p_0(U_{\mathrm a})
   =\operatorname{Tr}\!\left(\rho\,
            U_{\mathrm a}^\dagger\ket0\bra0 U_{\mathrm a}\right),
 \qquad
 p_0\bigl(R_y(-\pi/2)\bigr)
   =\frac{1+\operatorname{Tr}(\rho X)}2.
 \label{eq:platform-analysis-probability}
\end{equation}
The rotation $R_y(-\pi/2)$ maps $\ket+$ to $\ket0$, giving a zero
outcome with probability one, whereas $I/2$ remains unchanged and
still gives probability $1/2$. Similarly, $R_x(\pi/2)$ before
readout measures $(1+\operatorname{Tr}(\rho Y))/2$.
The pulse implementations of these rotations for each encoding
are developed in Section~\ref{sec:quantum-gates}. Repeated readout with
different analysis rotations can therefore determine the relative
phase information as well as the logical populations.
Testing a gate requires its action on independent input states;
the populations of one output cannot establish its action on an
arbitrary superposition. Appendices~\ref{sec:calibration-state-tomography}
and~\ref{sec:calibration-process-tomography} develop state and gate
reconstruction; Appendix~\ref{sec:calibration-gate-fidelity}
defines the quantitative comparison with a target gate.

